
\documentclass[11pt]{article}

\usepackage[utf8]{inputenc}
\usepackage[T1]{fontenc}
\usepackage[spanish,es-noshorthands,es-tabla]{babel}
\usepackage[a4paper,margin=2.1cm]{geometry}
\usepackage{amsmath,amssymb,amsfonts}
\usepackage{booktabs}
\usepackage{longtable}
\usepackage{tabularx}
\usepackage{array}
\usepackage{xcolor}
\usepackage{hyperref}
\usepackage{enumitem}
\usepackage{listings}
\usepackage{graphicx}
\usepackage{ragged2e}

\hypersetup{
  colorlinks=true,
  linkcolor=blue!55!black,
  urlcolor=blue!55!black,
  citecolor=blue!55!black
}

\newcommand{\code}[1]{\path{#1}}
\newcommand{\CaputoD}{\,{}^{C}D_t^q}
\newcommand{\ii}{\mathrm{i}}
\newcommand{\R}{\mathbb{R}}
\newcommand{\TARGET}{\texttt{TARGET}}

\setlength{\emergencystretch}{3em}
\setlength{\tabcolsep}{4pt}
\renewcommand{\arraystretch}{1.15}
\setlist[itemize]{leftmargin=1.2em}
\setlist[enumerate]{leftmargin=1.4em}

\newcolumntype{L}[1]{>{\RaggedRight\arraybackslash}p{#1}}
\DeclareRobustCommand{\codepath}[1]{\begingroup\Urlmuskip=0mu plus 2mu\path{#1}\endgroup}
\pdfstringdefDisableCommands{%
  \def\codepath#1{#1}%
  \def\code#1{#1}%
  \def\TARGET{TARGET}%
}


\lstdefinestyle{cmd}{
  basicstyle=\ttfamily\small,
  breaklines=true,
  columns=fullflexible,
  frame=single,
  rulecolor=\color{black!25},
  backgroundcolor=\color{black!3}
}
\lstset{style=cmd}

\title{Protocolo de pruebas para el caso de estudio\\
\large Chua fraccionario no suave, candidatos Lur'e/Machado y comparación con Danca}
\author{Hidden Attractors Fractional Order}
\date{\today}

\begin{document}
\maketitle
\tableofcontents

\section{Objetivo del documento}

Este documento describe cómo ejecutar, documentar e interpretar todas las pruebas disponibles en el repositorio para el caso de estudio del sistema de Chua fraccionario no suave. No sustituye el reporte matemático general; su función es convertir la teoría y los módulos de código en una ruta de trabajo reproducible.

El caso de referencia externa es el atractor oculto reportado por Danca en el sistema de Chua no suave de orden fraccionario. La comparación debe hacerse con cautela: Danca usa integración Caputo mediante Adams--Bashforth--Moulton con historia completa, mientras que la ruta principal del repositorio usa EFORK con memoria truncada \(L_m\). Por tanto, la comparación es geométrica, numérica y metodológica; no es identidad exacta de algoritmos.

\section{Estructura operativa del repositorio}

La versión de librería está en:

\begin{lstlisting}
version_2/
\end{lstlisting}

La estructura mínima que se usa en este protocolo es:

\begin{center}
\small
\begin{tabular}{@{}L{0.35\textwidth}L{0.56\textwidth}@{}}
\toprule
Ruta & Papel dentro del protocolo\\
\midrule
\codepath{hidden_attractors/} & Paquete importable. Incluye modelo, análisis, gráficas, clasificadores, workflows y backends nativos.\\
\codepath{examples/} & Ejemplos cortos para verificar instalación, equilibrios, candidatos y gráficas.\\
\codepath{tools/cli/} & Envolturas CLI mantenidas para workflows públicos.\\
\codepath{tools/legacy/} & Scripts de investigación todavía útiles para búsqueda Lur'e, búsqueda sesgada, Machado, Danca y corridas históricas. No todo es API pública.\\
\codepath{outputs/} & Corridas, CSV, JSON, trayectorias, figuras y candidatos ya generados.\\
\codepath{configs/} & Configuraciones YAML para corridas largas o búsquedas extendidas.\\
\bottomrule
\end{tabular}
\end{center}

Antes de ejecutar pruebas largas conviene instalar la librería en modo editable:

\begin{lstlisting}
cd version_2
python -m pip install -e .
\end{lstlisting}

Para pruebas de desarrollo:

\begin{lstlisting}
python -m pip install -e ".[dev]"
python -m pytest -q
\end{lstlisting}

Para habilitar análisis externos opcionales:

\begin{lstlisting}
python -m pip install -e ".[analysis]"
\end{lstlisting}

\section{Sistema de Chua fraccionario no suave}

El sistema estudiado es

\begin{align}
{}^C D_t^q x &= \alpha\left(y-x-f(x)\right),\\
{}^C D_t^q y &= x-y+z,\\
{}^C D_t^q z &= -\beta y-\gamma z,
\end{align}

con \(0<q\leq 1\), y no linealidad no suave

\begin{equation}
f(x)=m_1x+\frac{m_0-m_1}{2}\left(|x+1|-|x-1|\right).
\end{equation}

Equivalente en forma saturada,

\begin{equation}
f(x)=m_1x+(m_0-m_1)\operatorname{sat}(x),
\qquad
\operatorname{sat}(x)=
\begin{cases}
-1,&x<-1,\\
x,&|x|\leq 1,\\
1,&x>1.
\end{cases}
\end{equation}

La librería separa el campo vectorial del contrato fraccionario. El campo \(F(X)\) se evalúa en Python mediante funciones como \codepath{rhs_piecewise}; el orden \(q\), el paso \(h\), la memoria \(L_m\), el tiempo final \(T\) y el burn-in pertenecen al integrador o workflow.

\subsection{Parámetros por defecto de la librería}

En la parte pública de \codepath{version_2}, los parámetros por defecto del modelo no suave son

\begin{equation}
\alpha=8.4562,\qquad
\beta=12.0732,\qquad
\gamma=0.0052,\qquad
m_0=-0.1768,\qquad
m_1=-1.1468.
\end{equation}

Una forma directa de verlos es:

\begin{lstlisting}
python examples/quickstart_equilibria.py
\end{lstlisting}

O desde Python:

\begin{lstlisting}
from hidden_attractors import chua_piecewise_parameters

p = chua_piecewise_parameters()
print(p)
\end{lstlisting}

\section{Ver la gráfica del atractor}

\subsection{Usar una trayectoria ya generada}

La librería trabaja con trayectorias en formato CSV con columnas:

\begin{equation}
(t,x,y,z).
\end{equation}

Ejemplo:

\begin{lstlisting}
python examples/dynamical_analysis_gallery.py ^
  --trajectory-csv outputs/extended_search/machado_targeted_verification_lm10_20260515_182252/trajectories/branch_0_mu_4p00000_theta_0p00000_reference_attractor.csv ^
  --output-dir outputs/examples/chua_nonsmooth_gallery
\end{lstlisting}

En Linux/macOS:

\begin{lstlisting}
python examples/dynamical_analysis_gallery.py \
  --trajectory-csv outputs/extended_search/machado_targeted_verification_lm10_20260515_182252/trajectories/branch_0_mu_4p00000_theta_0p00000_reference_attractor.csv \
  --output-dir outputs/examples/chua_nonsmooth_gallery
\end{lstlisting}

Este ejemplo genera, típicamente:

\begin{itemize}
\item \codepath{phase_space_3d.png};
\item \codepath{phase_projections.png};
\item \codepath{time_series.png};
\item \codepath{bifurcation_diagram.png};
\item \codepath{summary.json};
\item \codepath{bifurcation_points.csv}.
\end{itemize}

\subsection{Código mínimo para graficar una trayectoria}

\begin{lstlisting}
from hidden_attractors.io import load_trajectory_csv
from hidden_attractors.plotting import (
    plot_phase_space,
    plot_phase_projections,
    plot_time_series,
)

traj = load_trajectory_csv("ruta/a/trayectoria.csv")

plot_phase_space(traj, "salidas/atractor_3d.png")
plot_phase_projections(traj, "salidas/proyecciones.png")
plot_time_series(traj, "salidas/series_temporales.png")
\end{lstlisting}

Lectura correcta: una gráfica de atractor muestra que la trayectoria es acotada y no trivial bajo el contrato numérico usado. No demuestra ocultedad.

\section{Puntos de equilibrio}

\subsection{Derivación algebraica}

En equilibrio se cumple

\begin{equation}
0=\alpha(y-x-f(x)),\qquad
0=x-y+z,\qquad
0=-\beta y-\gamma z.
\end{equation}

De la tercera ecuación,

\begin{equation}
z=-\frac{\beta}{\gamma}y,
\end{equation}

y de la segunda,

\begin{equation}
z=y-x.
\end{equation}

Entonces,

\begin{equation}
y-x=-\frac{\beta}{\gamma}y
\quad\Longrightarrow\quad
(\beta+\gamma)y=\gamma x
\quad\Longrightarrow\quad
y=\frac{\gamma}{\beta+\gamma}x.
\end{equation}

Además, de la primera ecuación,

\begin{equation}
y=x+f(x).
\end{equation}

Por tanto,

\begin{equation}
f(x)=y-x=\frac{\gamma}{\beta+\gamma}x-x
=
-\frac{\beta}{\beta+\gamma}x.
\end{equation}

Si \(x=0\), se obtiene

\begin{equation}
E_0=(0,0,0).
\end{equation}

En las ramas exteriores, para \(x>1\),

\begin{equation}
f(x)=m_1x+(m_0-m_1),
\end{equation}

y para \(x<-1\),

\begin{equation}
f(x)=m_1x-(m_0-m_1).
\end{equation}

Sea

\begin{equation}
s=-\frac{\beta}{\beta+\gamma}.
\end{equation}

La condición escalar exterior es

\begin{equation}
m_1x\pm(m_0-m_1)=sx.
\end{equation}

Así,

\begin{equation}
x_+=-\frac{m_0-m_1}{m_1-s},
\qquad
x_-=\frac{m_0-m_1}{m_1-s}.
\end{equation}

Finalmente,

\begin{equation}
E_\pm =
\left(
x_\pm,\,
x_\pm+f(x_\pm),\,
f(x_\pm)
\right).
\end{equation}

\subsection{Líneas para calcularlos con la librería}

\begin{lstlisting}
from hidden_attractors import chua_piecewise_parameters
from hidden_attractors.models import equilibria_piecewise, rhs_piecewise
import numpy as np

p = chua_piecewise_parameters()
eqs = equilibria_piecewise(p)

for name, point in eqs.items():
    residual = np.linalg.norm(rhs_piecewise(point, p))
    print(name, point, residual)
\end{lstlisting}

O directamente:

\begin{lstlisting}
python examples/quickstart_equilibria.py
\end{lstlisting}

\section{Jacobiano y estabilidad fraccionaria local}

La pendiente local de la no linealidad por tramos es

\begin{equation}
a(x)=
\begin{cases}
m_1,& |x|>1,\\
m_0,& |x|<1.
\end{cases}
\end{equation}

El Jacobiano por región es

\begin{equation}
J(a)=
\begin{bmatrix}
-\alpha(1+a)&\alpha&0\\
1&-1&1\\
0&-\beta&-\gamma
\end{bmatrix}.
\end{equation}

Para un sistema fraccionario conmensurado

\begin{equation}
{}^C D_t^q \xi=J(a)\xi,
\end{equation}

el criterio de Matignon exige

\begin{equation}
|\arg(\lambda_i)|>\frac{q\pi}{2}
\end{equation}

para estabilidad local asintótica. Si un equilibrio viola esa condición, es localmente inestable. Sin embargo, que un equilibrio sea inestable no implica que el atractor candidato sea autoexcitado; para eso se necesita verificar si su cuenca intersecta alguna vecindad de equilibrio.

\section{Forma Lur'e del sistema}

Escribimos

\begin{equation}
X=
\begin{bmatrix}
x\\y\\z
\end{bmatrix},
\qquad
\sigma=r^T X=x,
\qquad
r=
\begin{bmatrix}
1\\0\\0
\end{bmatrix}.
\end{equation}

Separando la parte lineal con pendiente exterior \(m_1\),

\begin{equation}
P=
\begin{bmatrix}
-\alpha(1+m_1)&\alpha&0\\
1&-1&1\\
0&-\beta&-\gamma
\end{bmatrix},
\qquad
b=
\begin{bmatrix}
-\alpha\\0\\0
\end{bmatrix}.
\end{equation}

La no linealidad escalar compatible es

\begin{equation}
\psi(\sigma)=(m_0-m_1)\operatorname{sat}(\sigma).
\end{equation}

Entonces

\begin{equation}
{}^C D_t^q X=PX+b\,\psi(r^TX).
\end{equation}

\section{Función de transferencia fraccionaria}

Para el bloque lineal fraccionario,

\begin{equation}
{}^C D_t^q X=PX+bu,\qquad \sigma=r^TX,
\end{equation}

la transferencia formal es

\begin{equation}
W_q(s)=r^T(s^q I-P)^{-1}b.
\end{equation}

Usando el parámetro espectral

\begin{equation}
\lambda=s^q,
\end{equation}

se define

\begin{equation}
\widehat W_q(\lambda)=r^T(\lambda I-P)^{-1}b.
\end{equation}

En frecuencia,

\begin{equation}
\lambda=(\ii\omega)^q
=
\omega^q
\left(
\cos\frac{q\pi}{2}
+
\ii\sin\frac{q\pi}{2}
\right).
\end{equation}

El código histórico usa tambien la convención

\begin{equation}
W_{\rm code}(\omega)=r^T(P-\lambda I)^{-1}b,
\end{equation}

por lo que

\begin{equation}
W_{\rm code}(\omega)=-\widehat W_q(\lambda).
\end{equation}

Esta diferencia de signo debe cuidarse al comparar ecuaciones del reporte con residuos emitidos por el código.

\section{Función descriptiva centrada clásica}

Para una entrada centrada

\begin{equation}
\sigma(t)=A\cos(\omega t),
\end{equation}

la función descriptiva clásica de la saturación escalada es

\begin{equation}
N(A)=\frac{Y_1}{A},
\end{equation}

donde \(Y_1\) es el primer armónico de \(\psi(A\cos\theta)\).

Para el caso no suave,

\begin{equation}
\psi(\sigma)=M\operatorname{sat}(\sigma),
\qquad
M=m_0-m_1.
\end{equation}

Si \(0<A\leq 1\),

\begin{equation}
N(A)=M.
\end{equation}

Si \(A>1\),

\begin{equation}
N(A)=\frac{2M}{\pi}
\left[
\arcsin\left(\frac{1}{A}\right)
+
\frac{\sqrt{A^2-1}}{A^2}
\right].
\end{equation}

La condición armónica, usando la convención de reporte, es

\begin{equation}
1+\widehat W_q((\ii\omega)^q)N(A)=0.
\end{equation}

Con la convención histórica del código,

\begin{equation}
1+W_{\rm code}(\omega)N(A)=0,
\end{equation}

pero recordando que \(W_{\rm code}=-\widehat W_q\). En el código, primero se resuelve

\begin{equation}
\Im W_{\rm code}(\omega)=0,
\end{equation}

y luego

\begin{equation}
k=-\frac{1}{\Re W_{\rm code}(\omega)}.
\end{equation}

La amplitud se obtiene resolviendo

\begin{equation}
N(A)=k.
\end{equation}

\subsection{Líneas para obtener semillas centradas}

Desde el flujo extendido:

\begin{lstlisting}
cd version_2
python tools/legacy/run_extended_search.py ^
  --config configs/chua_fractional_nonsmooth.yaml
\end{lstlisting}

En Linux/macOS:

\begin{lstlisting}
cd version_2
python tools/legacy/run_extended_search.py \
  --config configs/chua_fractional_nonsmooth.yaml
\end{lstlisting}

La parte centrada usa internamente:

\begin{itemize}
\item \codepath{find_omega_k_candidates};
\item \codepath{centered_candidates_from_nyquist};
\item \codepath{solve_amplitude_from_k};
\item \codepath{build_fractional_seed}.
\end{itemize}

Archivos esperados:

\begin{itemize}
\item \codepath{rho_H_diagnostics.csv};
\item \codepath{biased_df_candidates.csv};
\item \codepath{continuation_summary.csv};
\item \codepath{continuation_paths.csv};
\item \codepath{extended_search_report.md}.
\end{itemize}

Aunque el archivo \codepath{biased_df_candidates.csv} contiene la familia sesgada, el flujo extendido también construye candidatos centrados desde Nyquist antes de la continuación.

\section{Función descriptiva sesgada}

La función descriptiva sesgada usa

\begin{equation}
\sigma(t)=\sigma_0+A\cos(\omega t).
\end{equation}

La salida

\begin{equation}
y(t)=\psi(\sigma_0+A\cos(\omega t))
\end{equation}

se expande en Fourier:

\begin{equation}
y(t)=y_{\rm mean}+\sum_{k=1}^{K}
\left(a_k\cos(k\omega t)+b_k\sin(k\omega t)\right).
\end{equation}

El código define el coeficiente complejo

\begin{equation}
Y_k=a_k-\ii b_k.
\end{equation}

Por tanto,

\begin{equation}
N(A,\sigma_0)=\frac{Y_1}{A}.
\end{equation}

La ventaja frente a la función descriptiva centrada es que explora oscilaciones alrededor de una media distinta de cero. Esto es importante en sistemas con atractores desplazados respecto al origen.

\subsection{Condición armónica sesgada}

La condición usada para filtrar candidatos es

\begin{equation}
R(A,\sigma_0,\omega)
=
1+W_{\rm code}(\omega)N(A,\sigma_0),
\end{equation}

y se busca

\begin{equation}
|R(A,\sigma_0,\omega)|\ll 1.
\end{equation}

Además se calcula un diagnóstico de armónicos superiores:

\begin{equation}
\rho_H=
\frac{\sum_{k=2}^{K}|W_q(k\omega)|\,|Y_k|}
{|W_q(\omega)|\,|Y_1|+\varepsilon}.
\end{equation}

Interpretación:

\begin{itemize}
\item \(\rho_H\) pequeño: la aproximación de primer armónico es más razonable;
\item \(\rho_H\) grande: la función descriptiva pierde confiabilidad como generador de semilla;
\item \(\rho_H\) no prueba existencia de atractor ni ocultedad.
\end{itemize}

\subsection{Líneas para búsqueda sesgada}

El script principal es el mismo:

\begin{lstlisting}
python tools/legacy/run_extended_search.py ^
  --config configs/chua_fractional_nonsmooth.yaml
\end{lstlisting}

La búsqueda sesgada usa internamente:

\begin{itemize}
\item una malla en \((A,\sigma_0,\omega)\);
\item opcionalmente Latin Hypercube Sampling;
\item familias \codepath{classical_biased} y \codepath{machado_biased};
\item filtro por \codepath{residual_abs};
\item desempate por \(\rho_H\).
\end{itemize}

Archivos relevantes:

\begin{itemize}
\item \codepath{biased_df_all_evaluations.csv}: todas las evaluaciones;
\item \codepath{biased_df_candidates.csv}: mejores candidatos seleccionados;
\item \codepath{rho_H_diagnostics.csv}: diagnóstico armónico;
\item \codepath{plots/rho_H_vs_A.png};
\item \codepath{plots/rho_H_vs_omega.png};
\item \codepath{plots/harmonic_spectrum_candidate_*.png}.
\end{itemize}

\section{Función descriptiva tipo Machado}

La familia tipo Machado se usa como generador auxiliar de semillas. Si \(N(A)\) es la función descriptiva base, se define

\begin{equation}
N_{\mu}(A)=\exp\{\mu\operatorname{Log}N(A)\},
\qquad \mu>0.
\end{equation}

En el caso no suave centrado, si \(N(A)>0\), se usa la rama real:

\begin{equation}
N_{\mu}(A)=N(A)^\mu.
\end{equation}

En el caso sesgado, como \(N(A,\sigma_0)\) puede ser complejo, el código usa

\begin{equation}
N_{\mu}(A,\sigma_0)
=
\exp\left(\mu\operatorname{Log}_{\ell}N(A,\sigma_0)\right),
\end{equation}

donde \(\ell\) es la rama seleccionada del logaritmo complejo.

\subsection{Advertencia teórica}

El parámetro \(\mu\) no es el orden fraccionario del sistema. El sistema causal sigue teniendo orden \(q\). El parámetro \(\mu\) solo deforma la función descriptiva usada para generar semillas. Por tanto:

\begin{equation}
\mu\neq q.
\end{equation}

La ruta Machado no prueba ciclos exactos de Caputo. Produce candidatos que deben simularse y verificarse con cuencas.

\subsection{Líneas para activar Machado}

En la configuración YAML deben existir valores como:

\begin{lstlisting}
machado:
  mu_values: [0.25, 0.5, 0.75, 1.0, 1.25, 1.5, 2.0, 3.0, 4.0]
  branch: 0
  zero_eps: 1.0e-12

biased_search:
  families:
    - classical_biased
    - machado_biased
\end{lstlisting}

Después:

\begin{lstlisting}
python tools/legacy/run_extended_search.py ^
  --config configs/chua_fractional_nonsmooth.yaml
\end{lstlisting}

Para una verificación refinada de corridas Machado:

\begin{lstlisting}
python tools/legacy/run_extended_search.py ^
  --config configs/chua_fractional_nonsmooth.yaml ^
  --mode corrida1_refined_verification
\end{lstlisting}

\section{Otras rutas de generación que no conviene olvidar}

Además de la función descriptiva centrada, sesgada y Machado, el repositorio incluye o conserva estas rutas:

\begin{center}
\small
\begin{tabular}{@{}L{0.32\textwidth}L{0.56\textwidth}@{}}
\toprule
Ruta & Papel\\
\midrule
\(\rho_H\)-diagnostics & No genera una semilla nueva por sí solo, pero filtra la confiabilidad armónica de una semilla.\\
Seed cloud search & Simula nubes de condiciones iniciales evitando vecindades de equilibrio para detectar candidatos no triviales.\\
Continuación en \(\eta\) & Transporta una semilla desde una no linealidad suavizada o aproximada hasta la no linealidad original.\\
Continuación multiparamétrica & Cambia \(\eta\), \(q\), \(\sigma_0\) o \(\mu\) según la familia del candidato.\\
Danca ABM replication & Ruta independiente para reproducir el caso reportado por Danca con predictor--corrector ABM.\\
Refined basin classification & Reanaliza celdas desconocidas comparándolas con trayectorias target.\\
Robustness overlay & Cambia \(h,L_m,T\) para revisar persistencia geométrica.\\
Sphere controls & Lanza trayectorias desde esferas alrededor de los equilibrios.\\
External complexity tools & Usa \codepath{nolds} o \codepath{antropy} para entropías, dimensiones o indicadores de series temporales.\\
\bottomrule
\end{tabular}
\end{center}

\section{Modo sweep: barridos sistemáticos antes de elegir candidatos}

El modo \emph{sweep} no debe confundirse con una prueba final de ocultedad. Su papel es exploratorio: recorre una familia de parámetros, genera semillas o trayectorias candidatas, calcula residuos y diagnósticos, y deja archivos resumidos para decidir qué casos merecen continuación larga, cuencas y verificación alrededor de equilibrios.

La idea general es:
\begin{equation}
\mathcal P=\{p_1,p_2,\ldots,p_N\}
\quad\Longrightarrow\quad
\{\text{semilla}(p_i),\text{residuo}(p_i),\rho_H(p_i),\text{clase}(p_i)\}_{i=1}^N.
\end{equation}

En este repositorio existen varios barridos. Los más importantes para el Chua fraccionario no suave son:

\begin{center}
\small
\begin{tabular}{@{}L{0.27\textwidth}L{0.62\textwidth}@{}}
\toprule
Modo & Qué explora\\
\midrule
\codepath{q_sweep} & Barre el orden fraccionario \(q\). Para cada \(q\), calcula Nyquist/DF, semilla, continuación y atractor final ligero. Sirve para saber en qué órdenes fraccionarios el cierre armónico y la continuación son más prometedores.\\
\codepath{df_compare} & Compara la función descriptiva clásica contra la extensión tipo Machado para uno o varios valores de \(\mu\). Sirve para saber si Machado aporta semillas distintas de las clásicas.\\
\codepath{machado_sweep} & Barrido completo de ramas, valores de \(\mu\) y fases \(\theta\). Es más costoso y está pensado para exploración amplia del Chua no suave.\\
\codepath{machado_sweep_fast} & Versión reducida del barrido Machado. Usa una malla menor de \(\mu\) y \(\theta\), por lo que sirve como prueba previa antes de una corrida completa.\\
\codepath{positive_x_basin_sweep} & Barrido heredado sobre regiones de cuenca, útil para explorar qué zonas iniciales caen en clases target. Debe tratarse como herramienta de cartografía de cuenca, no como prueba independiente de ocultedad.\\
\bottomrule
\end{tabular}
\end{center}

\subsection{Barrido del orden fraccionario \texorpdfstring{\(q\)}{q}}

El modo \codepath{q_sweep} responde a la pregunta:

\begin{quote}
¿Para qué valores de \(q\) la ruta Nyquist/función descriptiva genera semillas continuables hacia un atractor no trivial?
\end{quote}

Con la entrada publica sin variables de entorno:

\begin{lstlisting}
hidden-attractors-unified-chua --run-mode q_sweep --q-values 0.97,0.98,0.985,0.99,0.995,1.0
\end{lstlisting}

Con malla uniforme:

\begin{lstlisting}
hidden-attractors-unified-chua --run-mode q_sweep --q-values linspace:0.96:1.0:17
\end{lstlisting}

La forma equivalente con modulo de Python es:

\begin{lstlisting}
python -m hidden_attractors.workflows.unified_chua --run-mode q_sweep --q-values 0.97,0.98,0.985,0.99,0.995,1.0
\end{lstlisting}

Las salidas esperadas son:

\begin{lstlisting}
q_order_sweep/
  q_sweep_summary.csv
  q_sweep_summary.json
  q_0p98500/
    fig01_nyquist_df_q_0p98500.pdf
    fig02_continuation_progress_q_0p98500.pdf
    fig03_final_attractor_q_0p98500.pdf
    continuation_summary_q_0p98500.json
    final_attractor_q_0p98500.csv
    final_attractor_q_0p98500.npz
    nyquist_df_samples_q_0p98500.csv
\end{lstlisting}

Lectura recomendada: elegir los valores de \(q\) donde el residuo armónico sea pequeño, la continuación no diverja y el atractor final sea acotado y no trivial. Después de esa preselección se ejecutan pruebas de robustez, cuencas y esferas alrededor de equilibrios. No se debe escribir ``atractor oculto'' únicamente a partir de \codepath{q_sweep}.

\subsection[Comparación clásica contra Machado: df\_compare]{Comparación clásica contra Machado: \codepath{df_compare}}

Este modo compara dos familias de semillas:

\begin{equation}
N_{\rm cl}(A),
\qquad
N_{\mu}(A)=N_{\rm cl}(A)^\mu.
\end{equation}

Se ejecuta, para cada familia, la cadena:

\begin{equation}
\text{Nyquist/DF}\rightarrow \text{semilla}\rightarrow \text{continuación}\rightarrow \text{verificación de ocultedad}.
\end{equation}

Comando básico:

\begin{lstlisting}
hidden-attractors-unified-chua --run-mode df_compare --q 0.99 --machado-mu 0.5
\end{lstlisting}

Para varios valores de \(\mu\):

\begin{lstlisting}
hidden-attractors-unified-chua --run-mode df_compare --q 0.99 --machado-mu-values 0.5,0.75,1.0,1.25,1.5,2.0,3.0
\end{lstlisting}

Salidas principales:

\begin{lstlisting}
df_seed_comparison/
  df_seed_comparison_summary.csv
  df_seed_comparison_summary.json
  classic/
  machado_mu_0p50000/
\end{lstlisting}

La comparación debe responder:
\begin{itemize}
\item si Machado produce amplitudes diferentes;
\item si la semilla resultante llega a una trayectoria acotada;
\item si la verificación preliminar detecta contacto con vecindades de equilibrio;
\item si la geometría final es distinta o equivalente a la semilla clásica.
\end{itemize}

\subsection{Barrido Machado completo}

El modo \codepath{machado_sweep} recorre ramas de Nyquist, órdenes descriptivos \(\mu\) y fases \(\theta\). Para cada combinación se construye

\begin{equation}
X_{\rm seed}^{(j,\mu,\theta)}
=
A_{j,\mu}\operatorname{Re}\left(v_j e^{i\theta}\right),
\end{equation}

donde \(j\) identifica la rama de Nyquist. La condición de compatibilidad usada en el barrido es

\begin{equation}
0<k_j<\Delta^\mu,
\qquad
\Delta=m_0-m_1.
\end{equation}

Comando completo:

\begin{lstlisting}
hidden-attractors-unified-chua --run-mode machado_sweep --model piecewise --output-dir runs_machado_sweep --q 0.99 --h 0.01 --memory-length 8 --t-transient 80 --t-keep 100
\end{lstlisting}

Salidas principales:

\begin{lstlisting}
machado_sweep/
  machado_sweep_summary.csv
  machado_sweep_summary.json
  branch_*/mu_*/theta_*/
\end{lstlisting}

Este barrido debe usarse cuando la búsqueda centrada/sesgada no produce candidatos suficientemente buenos o cuando se quiere justificar una exploración más amplia de fases.

\subsection{Barrido Machado rápido}

El modo \codepath{machado_sweep_fast} es la versión recomendada para una primera corrida. Usa una malla reducida:

\begin{align}
\text{rama 0: }&\mu\in\{0.75,1.00,1.25,1.50,2.00,3.00\},\\
\text{rama 1: }&\mu\in\{0.50,0.75,1.00,1.10,1.25,1.40\},\\
&\theta\in\left\{0,\frac{\pi}{3},\frac{2\pi}{3},\pi,\frac{4\pi}{3},\frac{5\pi}{3}\right\}.
\end{align}

Comando:

\begin{lstlisting}
hidden-attractors-unified-chua --run-mode machado_sweep_fast --model piecewise --output-dir runs_machado_sweep_fast --q 0.99 --h 0.01 --memory-length 8 --t-transient 80 --t-keep 100
\end{lstlisting}

Para una prueba mínima:

\begin{lstlisting}
hidden-attractors-unified-chua --run-mode machado_sweep_fast --machado-sweep-max-candidates 1
\end{lstlisting}

Salidas:

\begin{lstlisting}
machado_sweep_fast/
  machado_sweep_fast_summary.csv
  machado_sweep_fast_summary.json
  branch_*/mu_*/theta_*/
\end{lstlisting}

\subsection{Cómo filtrar un sweep}

Un sweep puede generar muchos candidatos. El filtrado recomendado es:

\begin{enumerate}
\item eliminar candidatos inválidos por incompatibilidad de ganancia;
\item ordenar por \codepath{residual_abs};
\item exigir \(\rho_H\) pequeño cuando esté disponible;
\item conservar únicamente trayectorias acotadas y no colapsadas;
\item descartar candidatos que converjan a equilibrio;
\item revisar contacto preliminar con vecindades de equilibrio;
\item pasar solo los mejores a continuación larga, robustez y cuencas.
\end{enumerate}

Operativamente, una tabla de sweep debe leerse así:

\begin{center}
\small
\begin{tabular}{@{}L{0.30\textwidth}L{0.58\textwidth}@{}}
\toprule
Columna & Interpretación\\
\midrule
\codepath{branch} & Rama de Nyquist usada para construir la semilla.\\
\codepath{mu} & Orden descriptivo auxiliar Machado. No es el orden de Caputo.\\
\codepath{theta} & Fase usada para colocar la semilla sobre la aproximación armónica.\\
\codepath{residual_abs} & Tamaño del error de cierre armónico \(|1+W_qN|\). Menor es mejor.\\
\codepath{rho_H} & Medida relativa de armónicos superiores. Menor indica mejor consistencia de primer armónico.\\
\codepath{survived} & Indica si la continuación o simulación preliminar no divergió.\\
\codepath{hiddenness_status} & Etiqueta operacional. No equivale por sí sola a prueba matemática.\\
\bottomrule
\end{tabular}
\end{center}

\subsection{Relación entre sweep y comparación con Danca}

Para comparar con Danca, el sweep sirve para localizar candidatos internos bajo el mismo sistema y un orden \(q\) cercano al reportado. La comparación debe separar tres cosas:

\begin{equation}
\begin{gathered}
\text{Danca/ABM con historia completa}\neq
\text{EFORK con memoria truncada}\\
\neq\text{semilla armónica tipo Weyl}.
\end{gathered}
\end{equation}

Por tanto, el reporte debe presentar:
\begin{itemize}
\item el atractor reportado/replicado por Danca;
\item el mejor candidato encontrado por \codepath{q_sweep}, \codepath{df_compare} o \codepath{machado_sweep_fast};
\item sus parámetros \(q,h,L_m,T,t_{\rm burn}\);
\item sus residuos armónicos y \(\rho_H\), si existen;
\item sus pruebas de cuenca y esfera alrededor de equilibrios;
\item una conclusión conservadora: compatible, descartado o pendiente de verificación, no ``demostrado'' salvo que todas las pruebas de cuenca lo respalden.
\end{itemize}


\section{Continuación numérica de un parámetro}

La continuación implementada en \codepath{multiparameter_continuation.py} usa una familia

\begin{equation}
{}^C D_t^q X=P X+b\,\psi_{\eta}(r^TX),
\qquad \eta\in[0,1].
\end{equation}

Para el modelo no suave, se usa una transición entre saturación suavizada y saturación exacta:

\begin{equation}
\psi_{\eta}(\sigma)
=
M\left[(1-\eta)\tanh\left(\frac{\sigma}{w}\right)
+
\eta\,\operatorname{sat}(\sigma)\right],
\qquad M=m_0-m_1.
\end{equation}

Casos:

\begin{equation}
\eta=0:\quad \psi_0(\sigma)=M\tanh(\sigma/w),
\end{equation}

\begin{equation}
\eta=1:\quad \psi_1(\sigma)=M\operatorname{sat}(\sigma).
\end{equation}

\subsection{Memoria en la continuación}

En sistemas de Caputo no basta pasar el último punto de una etapa a la siguiente. El código conserva una ventana discreta de historia:

\begin{equation}
\mathcal H_n=
\{(t_{n-M},X_{n-M}),\ldots,(t_n,X_n)\},
\qquad
M=\left\lceil\frac{L_m}{h}\right\rceil.
\end{equation}

Esta ventana se transporta mediante \codepath{FractionalHistory}. La continuación sigue siendo una heurística numérica, no una prueba de existencia exacta.

\subsection{Archivos de salida}

\begin{itemize}
\item \codepath{continuation_summary.csv}: una fila por paso de continuación;
\item \codepath{continuation_paths.csv}: ruta de estados finales;
\item columnas relevantes:
\begin{itemize}
\item \codepath{eta};
\item \codepath{q};
\item \codepath{chi_name};
\item \codepath{chi_value};
\item \codepath{bounded};
\item \codepath{diverged};
\item \codepath{equilibrium_hit};
\item \codepath{attractor_label};
\item \codepath{memory_carried};
\item \codepath{fft_dominant_frequency};
\item \codepath{psd_entropy}.
\end{itemize}
\end{itemize}

\section{Continuación multiparamétrica}

El cronograma de continuación se forma con tres posibles direcciones:

\begin{equation}
q_{\rm start}\to q_{\rm target},
\end{equation}

\begin{equation}
\chi_{\rm start}\to \chi_{\rm target},
\end{equation}

\begin{equation}
\eta_{\rm start}\to \eta_{\rm target}.
\end{equation}

La variable \(\chi\) depende de la familia del candidato:

\begin{equation}
\chi=
\begin{cases}
\mu,&\text{familia Machado},\\
\sigma_0,&\text{familia sesgada clásica}.
\end{cases}
\end{equation}

La configuración típica contiene:

\begin{lstlisting}
continuation:
  eta_start: 0.0
  eta_target: 1.0
  eta_steps: 6
  q_start: 0.9998
  q_target: 0.9998
  q_steps: 1
  chi_steps: 1
  h: 0.01
  memory_length: 10.0
  t_step: 60.0
  smooth_width: 0.2
  max_candidates: 10
\end{lstlisting}

\section{Candidatos finales y filtrado}

La librería pública carga tres candidatos finales mediante

\begin{lstlisting}
hidden-attractors-list-candidates
\end{lstlisting}

o

\begin{lstlisting}
python examples/list_final_candidates.py
\end{lstlisting}

Desde Python:

\begin{lstlisting}
from hidden_attractors import load_final_candidate_records

records = load_final_candidate_records()

for r in records:
    print(r.candidate_id, r.route, r.q, r.seed, r.robust_start)
\end{lstlisting}

\subsection{Criterios de filtrado recomendados}

Un candidato no debe seleccionarse solo por verse bien. El orden sugerido es:

\begin{enumerate}
\item Residuo armónico pequeño:
\[
|1+W_qN|\ll 1.
\]
\item \(\rho_H\) pequeño:
\[
\rho_H<\rho_{\rm max}.
\]
\item Semilla finita:
\[
\|X_0\|<\infty.
\]
\item Continuación sobreviviente:
\[
\texttt{bounded=True},\qquad \texttt{diverged=False}.
\]
\item No convergencia a equilibrio:
\[
\texttt{equilibrium\_hit=False}.
\]
\item Rango post-transitorio no trivial:
\[
\max\{\operatorname{range}_x,\operatorname{range}_y,\operatorname{range}_z\}>\varepsilon.
\]
\item Robustez bajo cambios de \(h,L_m,T\).
\item Ausencia de contacto target desde vecindades de equilibrios bajo radios probados.
\end{enumerate}

\section{Pruebas de ocultedad}

La definición operacional es:

\begin{quote}
Un atractor es oculto si su cuenca de atracción no intersecta ninguna vecindad abierta de los puntos de equilibrio. Si sí intersecta alguna vecindad de un equilibrio, es autoexcitado.
\end{quote}

En cómputo finito no se prueba una vecindad abierta completa. Se prueban radios y muestras:

\begin{equation}
X_0=E_i+r u_j,
\qquad
\|u_j\|=1,
\qquad
r\in\mathcal R.
\end{equation}

Si alguna trayectoria desde \(E_i+r u_j\) cae en la clase target, entonces hay evidencia contra la ocultedad bajo ese contrato.

\section{Cuencas de atracción}

Las cuencas se clasifican con etiquetas operativas:

\begin{center}
\small
\begin{tabular}{@{}cL{0.32\textwidth}L{0.46\textwidth}@{}}
\toprule
ID & Etiqueta & Interpretación\\
\midrule
0 & \codepath{equilibrium} & La trayectoria converge o queda cerca de un equilibrio.\\
1 & \codepath{target_positive} & La trayectoria cae en la clase target positiva.\\
2 & \codepath{target_negative} & La trayectoria cae en la clase target negativa.\\
3 & \codepath{infinity} & La trayectoria diverge o excede cota.\\
4 & \codepath{unknown} & La clasificación gruesa no decide.\\
5 & \codepath{numerical_failure} & Fallo numérico o datos no finitos.\\
6 & \codepath{bounded_other} & Clase refinada: acotada pero no suficientemente cercana al target.\\
\bottomrule
\end{tabular}
\end{center}

El corte de cuenca típico se hace en un plano, por ejemplo:

\begin{equation}
z=z_\ast,
\qquad
(x_0,y_0)\in[x_{\min},x_{\max}]\times[y_{\min},y_{\max}].
\end{equation}

Si quedan celdas \codepath{unknown}, se usa \codepath{refined_basin}.

\section{Prueba por esferas alrededor de equilibrios}

\subsection{Ejecución}

\begin{lstlisting}
hidden-attractors-sphere-controls --help
\end{lstlisting}

Ejecución con parámetros explícitos:

\begin{lstlisting}
hidden-attractors-sphere-controls ^
  --top-n 3 ^
  --radii 1e-5,3e-5,1e-4,3e-4,1e-3 ^
  --samples-per-radius 500 ^
  --chunks 6 ^
  --q 0.9998 ^
  --h 0.01 ^
  --memory-length 10.0 ^
  --t-final 1500.0 ^
  --t-burn 100.0
\end{lstlisting}

En Linux/macOS:

\begin{lstlisting}
hidden-attractors-sphere-controls \
  --top-n 3 \
  --radii 1e-5,3e-5,1e-4,3e-4,1e-3 \
  --samples-per-radius 500 \
  --chunks 6 \
  --q 0.9998 \
  --h 0.01 \
  --memory-length 10.0 \
  --t-final 1500.0 \
  --t-burn 100.0
\end{lstlisting}

\subsection{Archivos que produce}

\begin{itemize}
\item \codepath{sphere_plan.csv}: todas las condiciones iniciales sobre esferas;
\item \codepath{sphere_raw_chunk_*.csv}: resultados parciales por proceso;
\item \codepath{sphere_raw.csv}: todos los resultados;
\item \codepath{sphere_cumulative_summary.csv}: conteos acumulados;
\item \codepath{sphere_decision.csv}: decisión operacional por candidato;
\item \codepath{attractor_robustness_raw.csv}: clasificación de robustez desde el punto robusto;
\item \codepath{attractor_robustness_summary.csv};
\item \codepath{top3_sphere_robustness_summary.json}.
\end{itemize}

\subsection[Cómo leer sphere\_decision.csv]{Cómo leer \codepath{sphere_decision.csv}}

Columnas importantes:

\begin{itemize}
\item \codepath{candidate_id};
\item \codepath{total_sphere_trajectories};
\item \codepath{total_target_hits};
\item \codepath{smallest_radius_with_target_hit};
\item \codepath{hiddenness_status}.
\end{itemize}

Lectura:

\begin{equation}
\texttt{total\_target\_hits}>0
\quad\Longrightarrow\quad
\text{no se sostiene ocultedad bajo esas esferas}.
\end{equation}

\begin{equation}
\texttt{total\_target\_hits}=0
\quad\Longrightarrow\quad
\text{compatible con ocultedad bajo radios y muestras probados}.
\end{equation}

No debe escribirse ``atractor oculto demostrado'' solo con esta condición.

\section{Pruebas de robustez}

\subsection{Robustness overlay}

Este workflow no prueba ocultedad. Compara la geometría del atractor al modificar:

\begin{equation}
h,\qquad L_m,\qquad T.
\end{equation}

Ejecución:

\begin{lstlisting}
hidden-attractors-robustness-overlay --help
\end{lstlisting}

Ejemplo:

\begin{lstlisting}
hidden-attractors-robustness-overlay ^
  --q 0.9998 ^
  --divergence-norm 120.0 ^
  --equilibrium-tol 1e-3 ^
  --max-store-points 6000 ^
  --max-metric-points 1000 ^
  --max-section-points 300
\end{lstlisting}

\subsection{Casos de robustez típicos}

\begin{center}
\small
\begin{tabular}{@{}L{0.22\textwidth}cccL{0.26\textwidth}@{}}
\toprule
Caso & \(h\) & \(L_m\) & \(T\) & Lectura\\
\midrule
\codepath{R0_base} & 0.01 & 10 & 1500 & Contrato base.\\
\codepath{R1_h_finer} & 0.005 & 10 & 1500 & Paso más fino.\\
\codepath{R2_h_coarser} & 0.02 & 10 & 1500 & Paso más grueso.\\
\codepath{R3_Lm_lower} & 0.01 & 5 & 1500 & Menor memoria.\\
\codepath{R4_Lm_higher} & 0.01 & 20 & 1500 & Mayor memoria.\\
\codepath{R5_t_longer} & 0.01 & 10 & 3000 & Tiempo más largo.\\
\bottomrule
\end{tabular}
\end{center}

\subsection{Archivos de salida}

\begin{itemize}
\item \codepath{robustness_overlay_config.json};
\item \codepath{metrics_<candidate>.csv};
\item \codepath{robustness_overlay_metrics.csv};
\item \codepath{robustness_overlay_summary.json};
\item \codepath{plots/overlay_<candidate>.png}.
\end{itemize}

\subsection{Cómo decidir si un candidato es robusto}

Un candidato es robusto bajo estos contratos si:

\begin{equation}
\texttt{bounded=True},
\qquad
\texttt{diverged=False},
\qquad
\texttt{equilibrium\_like=False},
\end{equation}

y las distancias relativas son pequeñas:

\begin{equation}
\texttt{range\_relative\_distance}\ll 1,
\end{equation}

\begin{equation}
\texttt{cloud\_median\_distance\_norm}\ll 1,
\end{equation}

\begin{equation}
\texttt{section\_median\_distance\_norm}\ll 1.
\end{equation}

\section{Refinamiento de cuencas desconocidas}

Si una cuenca tiene muchas celdas \codepath{unknown}, se usa:

\begin{lstlisting}
hidden-attractors-refined-basin --help
\end{lstlisting}

Ejemplo:

\begin{lstlisting}
hidden-attractors-refined-basin ^
  --source-dir outputs/basin_compare_danca_project_h001_grid101_20260517 ^
  --chunks 8 ^
  --tail-fraction-start 0.5 ^
  --max-cloud-points 700 ^
  --max-section-points 250 ^
  --max-score 0.95
\end{lstlisting}

El método reintegra cada celda desconocida y compara contra dos referencias:

\begin{equation}
\mathcal A_+,\qquad \mathcal A_-.
\end{equation}

Calcula:

\begin{itemize}
\item distancia entre nubes de cola;
\item distancia entre secciones de Poincaré;
\item distancia relativa de rangos;
\item diferencia relativa de frecuencia FFT dominante;
\item score ponderado.
\end{itemize}

La decisión se expresa como:

\begin{equation}
\texttt{matched\_target\_reference}
\end{equation}

o

\begin{equation}
\texttt{bounded\_noncollapsed\_reference\_mismatch}.
\end{equation}

\section{Análisis adicionales implementados}

\subsection{FFT y entropía espectral}

Para una señal \(x_n\), se resta la media y se usa una ventana de Hann:

\begin{equation}
\tilde x_n=(x_n-\bar x)w_n.
\end{equation}

El espectro de potencia discreto es

\begin{equation}
S_k=|\operatorname{rFFT}(\tilde x)_k|^2.
\end{equation}

La frecuencia dominante es

\begin{equation}
f_{\max}=f_{\arg\max_{k\geq 1}S_k}.
\end{equation}

La entropía espectral normalizada se calcula como

\begin{equation}
H_{\rm PSD}
=
-\frac{1}{\log N}
\sum_k p_k\log(p_k+\varepsilon),
\qquad
p_k=\frac{S_k}{\sum_j S_j}.
\end{equation}

\subsection{Sección de Poincaré operacional}

El código usa cruces hacia arriba del plano \(x=0\):

\begin{equation}
x_{n-1}<0\leq x_n,
\qquad
\dot x_n>0.
\end{equation}

En un sistema fraccionario no se interpreta esta sección como mapa de Poincaré exacto con propiedad de semigrupo. Se usa como nube geométrica para comparar atractores.

\subsection{Distancia entre nubes}

Dadas dos nubes \(A\) y \(B\), se define una distancia simétrica basada en vecinos más cercanos:

\begin{equation}
d(A,B)=
\operatorname{median}
\left(
\{\min_{b\in B}\|a-b\|:a\in A\}
\cup
\{\min_{a\in A}\|b-a\|:b\in B\}
\right).
\end{equation}

\subsection{Bifurcación por postprocesamiento}

La librería extrae máximos, mínimos o muestras de trayectorias ya calculadas:

\begin{lstlisting}
from hidden_attractors.analysis import bifurcation_points_from_trajectories
from hidden_attractors.plotting import plot_bifurcation_diagram
\end{lstlisting}

Esto no es continuación numérica formal. Sirve para diagramas exploratorios.

\subsection{Medidas externas de complejidad}

Si se instala \codepath{nolds} o \codepath{antropy}:

\begin{lstlisting}
from hidden_attractors.integrations import compute_complexity_measures

metrics = compute_complexity_measures(traj[:, 1], backend="auto")
print(metrics)
\end{lstlisting}

Estas medidas pueden incluir entropía de permutación, entropía muestral, dimensión de correlación, DFA o estimadores de Lyapunov de serie temporal. No son prueba de ocultedad.

\section{Cómo elegir el mejor candidato}

Un ranking razonable puede organizarse así:

\begin{center}
\small
\begin{tabular}{@{}cL{0.42\textwidth}L{0.36\textwidth}@{}}
\toprule
Prioridad & Criterio & Archivo\\
\midrule
1 & Residuo armónico pequeño & \codepath{biased_df_candidates.csv}, \codepath{rho_H_diagnostics.csv}\\
2 & Armónicos superiores atenuados & \codepath{rho_H_diagnostics.csv}\\
3 & Supervivencia por continuación & \codepath{continuation_summary.csv}\\
4 & No convergencia a equilibrio & \codepath{continuation_summary.csv}\\
5 & Trayectoria acotada no trivial & \codepath{metrics_*.csv}, \codepath{trajectory_metrics}\\
6 & Robustez ante \(h,L_m,T\) & \codepath{robustness_overlay_metrics.csv}\\
7 & No contacto desde esferas & \codepath{sphere_decision.csv}\\
8 & Cuenca refinada sin contacto local & \codepath{refined_basin_summary.json}\\
9 & Comparación geométrica contra Danca & \codepath{danca2017_vs_project_candidates_report.md}\\
\bottomrule
\end{tabular}
\end{center}

Una regla conservadora:

\begin{equation}
\text{mejor candidato}
=
\arg\min
\left[
w_1|R|+
w_2\rho_H+
w_3d_{\rm cloud}+
w_4d_{\rm section}
+
w_5 I_{\rm contact}
\right],
\end{equation}

donde \(I_{\rm contact}=+\infty\) si hay contacto target desde vecindades de equilibrio. En términos prácticos: cualquier contacto target desde una vecindad de equilibrio pesa más que una buena gráfica 3D.

\section{Gráficas recomendadas para cada candidato}

Para cada candidato final debe generarse:

\begin{enumerate}
\item atractor 3D;
\item proyecciones \(xy,xz,yz\);
\item series temporales;
\item espectro FFT o PSD;
\item sección \(x=0\);
\item overlay de robustez;
\item cuenca \(xy\);
\item cuencas \(xz\) y \(yz\), si están disponibles;
\item zoom alrededor de \(E_0,E_+,E_-\);
\item esferas alrededor de \(E_0,E_+,E_-\);
\item target hits contra radio;
\item tabla final de contratos numéricos.
\end{enumerate}

\section{Comparación con el atractor reportado por Danca}

\subsection{Ruta Danca}

El script específico es:

\begin{lstlisting}
python tools/legacy/danca2017_chua_abm_replication.py --help
\end{lstlisting}

Ejecución rápida no científica:

\begin{lstlisting}
python tools/legacy/danca2017_chua_abm_replication.py ^
  --job all ^
  --quick ^
  --workers 4 ^
  --output-dir outputs/danca_quick_test
\end{lstlisting}

Ejecución más seria:

\begin{lstlisting}
python tools/legacy/danca2017_chua_abm_replication.py ^
  --job all ^
  --workers 4 ^
  --q 0.9998 ^
  --h 0.05 ^
  --t-final 500.0 ^
  --transient 250.0 ^
  --delta 0.01 ^
  --local-samples-per-unstable-eq 100 ^
  --figure-local-trajectories 80 ^
  --output-dir outputs/danca2017_chua_abm_reproduction
\end{lstlisting}

\subsection{Qué produce}

\begin{itemize}
\item \codepath{run_config.json};
\item \codepath{danca_reference_summary.json};
\item \codepath{danca_hiddenness_decision.json};
\item \codepath{fig03_danca2017_chua_abm_replica.png};
\item \codepath{fig03_danca2017_chua_abm_replica.pdf};
\item \codepath{project_best_two_figure_summary.csv};
\item \codepath{danca2017_vs_project_candidates_report.md}.
\end{itemize}

\subsection{Diferencias metodológicas}

\begin{center}
\small
\begin{tabular}{@{}L{0.28\textwidth}L{0.29\textwidth}L{0.31\textwidth}@{}}
\toprule
Aspecto & Danca & Repositorio\\
\midrule
Integrador & ABM predictor--corrector & EFORK con memoria truncada; ABM solo en script de réplica\\
Memoria & Historia completa de Caputo & Ventana finita \(L_m\)\\
Localización & Búsqueda por prueba y error reportada & Semillas por Lur'e, DF centrada, DF sesgada, Machado y nubes\\
Ocultedad & Vecindades de equilibrios con \(\delta=0.01\) & Esferas, cuencas, refinamiento y comparación target\\
Salida & Figura y argumento del artículo & CSV/JSON/figuras reproducibles\\
Lectura & Atractor oculto reportado & Candidato compatible/no compatible bajo contrato\\
\bottomrule
\end{tabular}
\end{center}

\subsection{Comparación mínima que debe aparecer en el reporte}

Para comparar un candidato del repositorio con Danca, incluir una tabla:

\begin{center}
\small
\begin{tabular}{@{}L{0.24\textwidth}L{0.27\textwidth}L{0.36\textwidth}@{}}
\toprule
Bloque & Dato & Interpretación\\
\midrule
Modelo & mismos \(\alpha,\beta,\gamma,m_0,m_1,q\) & Comparación válida a nivel de sistema\\
Integrador & ABM vs EFORK & No son numéricamente idénticos\\
Condición inicial & semilla Danca localizada o no publicada & Si no está publicada, declararlo\\
Atrayente observado & rangos, FFT, nube, sección & Comparación geométrica\\
Equilibrios & \(E_0,E_+,E_-\) & Base de prueba de ocultedad\\
Pruebas locales & \(\delta\) Danca vs radios \(\mathcal R\) & Comparar tamaños de vecindad\\
Resultado & target/no target desde vecindades & Decide compatibilidad con ocultedad\\
\bottomrule
\end{tabular}
\end{center}

\section{Flujo completo recomendado}

\begin{enumerate}
\item Instalar la librería:
\begin{lstlisting}
cd version_2
python -m pip install -e ".[dev]"
\end{lstlisting}

\item Verificar equilibrios:
\begin{lstlisting}
python examples/quickstart_equilibria.py
\end{lstlisting}

\item Graficar una trayectoria existente:
\begin{lstlisting}
python examples/dynamical_analysis_gallery.py ^
  --trajectory-csv ruta/a/trayectoria.csv ^
  --output-dir outputs/examples/galeria_prueba
\end{lstlisting}

\item Ejecutar búsqueda extendida Lur'e/DF/Machado:
\begin{lstlisting}
python tools/legacy/run_extended_search.py ^
  --config configs/chua_fractional_nonsmooth.yaml
\end{lstlisting}

\item Revisar candidatos:
\begin{lstlisting}
python examples/list_final_candidates.py
hidden-attractors-list-candidates
\end{lstlisting}

\item Ejecutar robustez:
\begin{lstlisting}
hidden-attractors-robustness-overlay
\end{lstlisting}

\item Ejecutar esferas de equilibrio:
\begin{lstlisting}
hidden-attractors-sphere-controls ^
  --top-n 3 ^
  --samples-per-radius 500 ^
  --chunks 6
\end{lstlisting}

\item Refinar cuenca si existen celdas \codepath{unknown}:
\begin{lstlisting}
hidden-attractors-refined-basin ^
  --source-dir outputs/basin_compare_danca_project_h001_grid101_20260517 ^
  --chunks 8
\end{lstlisting}

\item Reproducir Danca:
\begin{lstlisting}
python tools/legacy/danca2017_chua_abm_replication.py ^
  --job all ^
  --workers 4 ^
  --output-dir outputs/danca2017_chua_abm_reproduction
\end{lstlisting}

\item Comparar Danca contra candidatos del proyecto:
\begin{lstlisting}
python tools/legacy/danca2017_chua_abm_replication.py ^
  --job report ^
  --output-dir outputs/danca2017_chua_abm_reproduction
\end{lstlisting}
\end{enumerate}

\section{Criterio final de redacción científica}

La conclusión debe redactarse con niveles de evidencia:

\begin{itemize}
\item ``La función descriptiva predice una semilla'' si solo se resolvió la condición armónica.
\item ``La continuación produce una trayectoria acotada no trivial'' si sobrevivió la homotopía.
\item ``El candidato es robusto bajo los contratos probados'' si persiste al cambiar \(h,L_m,T\).
\item ``El candidato es compatible con ocultedad bajo los radios probados'' si no hay target hits desde esferas.
\item ``El candidato no está soportado como oculto bajo este contrato'' si hay target hits desde alguna vecindad de equilibrio.
\item ``Atractor oculto verificado numéricamente'' solo si se documentan todos los equilibrios, radios, cuencas, robustez, integración y ausencia de intersección de cuenca con vecindades probadas. Aun así, debe decirse que es verificación numérica finita.
\end{itemize}

\section{Checklist para una corrida defendible}

\begin{enumerate}
\item ¿Se fijaron \(\alpha,\beta,\gamma,m_0,m_1,q\)?
\item ¿Se reportó el integrador?
\item ¿Se reportaron \(h,L_m,T,t_{\rm burn}\)?
\item ¿Se calcularon \(E_0,E_+,E_-\)?
\item ¿Se revisó Matignon localmente?
\item ¿Se documentó cómo se obtuvo la semilla?
\item ¿Se guardó \(|1+W_qN|\)?
\item ¿Se guardó \(\rho_H\)?
\item ¿Se realizó continuación con memoria?
\item ¿Se probó robustez?
\item ¿Se generaron cuencas?
\item ¿Se probaron esferas alrededor de cada equilibrio?
\item ¿Se comparó contra Danca con el mismo sistema y explicando diferencias de integrador?
\item ¿Se evitó declarar ocultedad solo por una gráfica?
\end{enumerate}


\section{Balance armónico implementado en el repositorio}
\label{sec:balance_armonico_detallado}

El balance armónico es la etapa que convierte la forma de Lur'e en una condición algebraica aproximada para proponer semillas iniciales. No es una prueba de existencia de ciclos límite exactos de Caputo. En el repositorio se usa como aproximación estacionaria tipo Weyl para construir una condición inicial que después se valida con integración causal de Caputo/EFORK o ABM.

Partimos de
\begin{equation}
\CaputoD X=PX+b\psi(r^TX),
\qquad \sigma=r^TX.
\end{equation}
Para un régimen dominante supuesto,
\begin{equation}
\sigma(t)\approx \sigma_0+A\cos(\omega t),
\end{equation}
la no linealidad se expande como
\begin{equation}
\psi(\sigma_0+A\cos \theta)
=Y_0+\sum_{k=1}^{\infty}
\left(a_k\cos(k\theta)+b_k\sin(k\theta)\right),
\qquad \theta=\omega t.
\end{equation}
El código usa coeficientes complejos
\begin{equation}
Y_k=a_k-\ii b_k.
\end{equation}
La función descriptiva de primer armónico es
\begin{equation}
N(A,\sigma_0)=\frac{Y_1}{A}.
\end{equation}
En el caso centrado, \(\sigma_0=0\), se escribe simplemente \(N(A)\).

La evaluación frecuencial del bloque lineal fraccionario no se hace con \(s=\ii\omega\) directamente, sino con
\begin{equation}
\lambda=(\ii\omega)^q=\omega^q\exp\left(\ii\frac{q\pi}{2}\right).
\end{equation}
Con la convención de reporte,
\begin{equation}
\widehat W_q(\lambda)=r^T(\lambda I-P)^{-1}b.
\end{equation}
El cierre de primer armónico queda
\begin{equation}
1+\widehat W_q((\ii\omega)^q)N(A,\sigma_0)=0.
\end{equation}
La convención histórica del código define
\begin{equation}
W_{\rm code}(\omega)=r^T(P-(\ii\omega)^qI)^{-1}b
=-\widehat W_q((\ii\omega)^q).
\end{equation}
Por eso, en los CSV históricos el residuo aparece como
\begin{equation}
R(A,\sigma_0,\omega)=1+W_{\rm code}(\omega)N(A,\sigma_0).
\end{equation}
Al interpretar resultados del código debe respetarse esa convención de signo.

\subsection{Condición de Nyquist usada para semillas centradas}

Para la función descriptiva centrada se busca primero una frecuencia \(\omega_0\) tal que
\begin{equation}
\Im W_{\rm code}(\omega_0)=0.
\end{equation}
Después se calcula
\begin{equation}
k=-\frac{1}{\Re W_{\rm code}(\omega_0)}.
\end{equation}
La amplitud se obtiene de
\begin{equation}
N(A)=k.
\end{equation}
Si \(k\) no pertenece al rango admisible de la función descriptiva, la rama se descarta. Para la saturación no suave con \(M=m_0-m_1>0\), el código exige
\begin{equation}
0<k\leq M.
\end{equation}

\subsection{Diagnóstico de armónicos superiores}

La condición armónica sólo conserva el primer armónico. Para medir la severidad de esa aproximación se calcula
\begin{equation}
\rho_H=
\frac{\displaystyle\sum_{k=2}^{K}|W_q(k\omega)|\,|Y_k|}
{|W_q(\omega)|\,|Y_1|+\varepsilon}.
\end{equation}
La lectura correcta es:
\begin{itemize}
\item \(\rho_H\) pequeño: la semilla armónica es más razonable como punto inicial;
\item \(\rho_H\) grande: los armónicos superiores no están suficientemente atenuuados;
\item \(\rho_H\) no prueba caos, ocultedad ni existencia de ciclo periódico exacto.
\end{itemize}

En el flujo extendido, los campos importantes son:
\begin{center}
\begin{tabular}{@{}ll@{}}
\toprule
Campo & Significado\\
\midrule
\codepath{residual_abs} & \(|1+W_{\rm code}N|\).\\
\codepath{rho_H} & Razón de armónicos superiores.\\
\codepath{harmonic_energy_ratio} & Energía relativa fuera del primer armónico.\\
\codepath{accepted_by_rhoH} & Filtro lógico comparado contra el umbral del YAML.\\
\bottomrule
\end{tabular}
\end{center}

\section{Cálculo explícito de semillas iniciales}
\label{sec:calculo_semillas_detallado}

El repositorio contiene tres mecanismos principales de semilla: centrada clásica, sesgada y Machado. Además conserva búsqueda por nubes de semillas externas a vecindades de equilibrio.

\subsection{Semilla centrada clásica}

Una vez obtenidos \((\omega_0,k,A)\), se introduce el bloque lineal modificado
\begin{equation}
P_0=P+kb r^T.
\end{equation}
La semilla armónica se construye con el autovector asociado a
\begin{equation}
\lambda_0=(\ii\omega_0)^q.
\end{equation}
El problema que se usa es
\begin{equation}
(P_0-\lambda_0 I)v\approx 0,
\qquad r^Tv=1.
\end{equation}
Luego, para una fase \(\theta\),
\begin{equation}
X_{\rm seed}(\theta)=A\Re\{v e^{\ii\theta}\}.
\end{equation}
En código, esta ruta corresponde a:
\begin{itemize}
\item \codepath{find_omega_k_candidates};
\item \codepath{solve_amplitude_from_k};
\item \codepath{build_fractional_seed}.
\end{itemize}

\subsection{Semilla sesgada}

Para una semilla sesgada se usa
\begin{equation}
\sigma(t)=\sigma_0+A\cos(\omega t).
\end{equation}
La componente media de la no linealidad, \(Y_0=y_{\rm mean}\), produce un equilibrio medio aproximado \(\bar X\). El código lo obtiene por mínimos cuadrados imponiendo simultáneamente
\begin{equation}
P\bar X+bY_0=0,
\qquad r^T\bar X=\sigma_0.
\end{equation}
La componente fundamental compleja \(V\) se obtiene resolviendo
\begin{equation}
(\lambda I-P)V=bY_1,
\qquad r^TV=A.
\end{equation}
La semilla sesgada queda
\begin{equation}
X_{\rm seed}(\theta)=\bar X+\Re\{Ve^{\ii\theta}\}.
\end{equation}
Esta ruta evita asumir que el atractor oscila alrededor del origen. Es especialmente útil cuando el candidato se localiza en una región desplazada de fase.

En código, esta construcción corresponde a:
\begin{itemize}
\item \codepath{fourier_coefficients_psi};
\item \codepath{N_biased};
\item \codepath{reconstruct_biased_lure_seed};
\item \codepath{search_biased_candidates}.
\end{itemize}

\subsection{Semilla Machado centrada}

Para el caso centrado no suave, la extensión tipo Machado usa
\begin{equation}
N_{\mu}(A)=N(A)^\mu,
\qquad \mu>0.
\end{equation}
Para cada rama de Nyquist se conserva la misma \(\omega_j\) y la misma ganancia lineal \(k_j\), pero la amplitud se recalcula con
\begin{equation}
N{(A_\mu)}^\mu=k_j,
\qquad\text{o equivalentemente}\qquad
N(A_\mu)=k_j^{1/\mu}.
\end{equation}
La compatibilidad exigida es
\begin{equation}
0<k_j<M^{\mu}.
\end{equation}
Después se usa el mismo autovector fraccionario:
\begin{equation}
X_{{\rm seed},\mu,j}(\theta)=A_{\mu,j}\Re\{v_j e^{\ii\theta}\}.
\end{equation}

\subsection{Semilla Machado sesgada}

Cuando \(N(A,\sigma_0)\) es complejo, se usa
\begin{equation}
N_{\mu}(A,\sigma_0)=
\exp\left(\mu\operatorname{Log}_{\ell}N(A,\sigma_0)\right),
\end{equation}
donde \(\ell\) fija la rama del logaritmo complejo. Esta ruta es más delicada porque puede introducir discontinuidades de rama. Por eso el código descarta valores cercanos a cero con \codepath{zero_eps}.

\subsection{Nubes de semillas}

La búsqueda por nubes no usa función descriptiva. Genera condiciones iniciales en una caja
\begin{equation}
[x_{\min},x_{\max}]\times[y_{\min},y_{\max}]\times[z_{\min},z_{\max}],
\end{equation}
mezclando malla cartesiana y Latin Hypercube Sampling. Después elimina puntos demasiado cercanos a los equilibrios:
\begin{equation}
\min_i\|X_0-E_i\|>\delta_{\rm eq}.
\end{equation}
La clasificación preliminar puede producir etiquetas como \codepath{candidate_hidden_like}, pero esa etiqueta sólo significa que la trayectoria observada fue acotada, no trivial y no pasó cerca de los equilibrios bajo el contrato corto probado. No equivale a \codepath{hidden_verified}.

\section{Criterios de robustez: pruebas implementadas y lectura}
\label{sec:robustez_todas_pruebas}

El repositorio implementa varias familias de robustez. Conviene distinguirlas porque no responden la misma pregunta.

\subsection{Robustez por superposición de trayectorias}

El workflow público \codepath{robustness_overlay} compara la geometría del mismo candidato cuando se cambia el contrato numérico. La lista estándar es:
\begin{center}
\begin{tabular}{@{}lllll@{}}
\toprule
Caso & \(h\) & \(L_m\) & \(T\) & Lectura\\
\midrule
\codepath{R0_base} & 0.01 & 10 & 1500 & Referencia.\\
\codepath{R1_h_finer} & 0.005 & 10 & 1500 & Sensibilidad a paso más fino.\\
\codepath{R2_h_coarser} & 0.02 & 10 & 1500 & Sensibilidad a paso más grueso.\\
\codepath{R3_Lm_lower} & 0.01 & 5 & 1500 & Sensibilidad a memoria menor.\\
\codepath{R4_Lm_higher} & 0.01 & 20 & 1500 & Sensibilidad a memoria mayor.\\
\codepath{R5_t_longer} & 0.01 & 10 & 3000 & Persistencia en tiempo largo.\\
\bottomrule
\end{tabular}
\end{center}

Para cada caso se guardan métricas:
\begin{itemize}
\item \codepath{bounded}, \codepath{diverged}, \codepath{equilibrium_like};
\item rangos \codepath{range_x}, \codepath{range_y}, \codepath{range_z};
\item varianzas de cola;
\item frecuencia FFT dominante;
\item entropía espectral;
\item número de puntos de sección;
\item distancia relativa de rangos frente a \codepath{R0_base};
\item diferencia relativa de FFT;
\item distancia mediana entre nubes de cola;
\item distancia mediana entre secciones de Poincaré.
\end{itemize}

Un candidato es robusto sólo en sentido operativo si, al menos:
\begin{equation}
\texttt{bounded=True},\qquad
\texttt{diverged=False},\qquad
\texttt{equilibrium\_like=False},
\end{equation}
para todos los casos probados, y además las distancias geométricas respecto de la referencia no son desproporcionadas. El código no convierte automáticamente esta comparación en prueba de ocultedad.

\subsection{Robustez por clasificador de cuenca desde el punto candidato}

El workflow \codepath{sphere_controls} incluye una prueba adicional desde el punto final robusto de cada candidato. Usa cuatro contratos:
\begin{center}
\begin{tabular}{@{}llll@{}}
\toprule
Caso & Cambio & Objetivo & Salida\\
\midrule
\codepath{R0_baseline} & contrato base & clase target inicial & \codepath{attractor_robustness_raw.csv}\\
\codepath{R1_longer} & aumenta \(T\) & persistencia temporal & \codepath{attractor_robustness_summary.csv}\\
\codepath{R2_shorter_memory} & reduce \(L_m\) & sensibilidad a memoria & \codepath{attractor_robustness_summary.csv}\\
\codepath{R3_refined_h} & reduce \(h\) & sensibilidad a paso & \codepath{attractor_robustness_summary.csv}\\
\bottomrule
\end{tabular}
\end{center}
El campo \codepath{robust_attractor=True} significa que el candidato cayó en clase target bajo todos esos contratos. No significa que sea oculto.

\subsection{Robustez por esferas alrededor de equilibrios}

La prueba decisiva contra ocultedad es local alrededor de cada equilibrio. Para cada equilibrio \(E_i\), cada radio \(r_j\) y cada dirección unitaria \(u_k\), se integra desde
\begin{equation}
X_0=E_i+r_j u_k.
\end{equation}
Si alguna trayectoria se clasifica como \codepath{target_positive} o \codepath{target_negative}, entonces existe evidencia contra la ocultedad bajo ese contrato:
\begin{equation}
\exists i,j,k:
\quad X_0=E_i+r_ju_k\longrightarrow \TARGET.
\end{equation}
La lectura del archivo \codepath{sphere_decision.csv} es:
\begin{itemize}
\item \codepath{not_supported_by_sphere_equilibrium_test}: hubo contacto target desde alguna vecindad de equilibrio;
\item \codepath{compatible_with_hiddenness_under_tested_spheres}: no se observó contacto target bajo los radios y muestras probados.
\end{itemize}
La segunda etiqueta no es demostración formal; sólo resume una prueba finita.

\subsection{Robustez por refinamiento de cuencas desconocidas}

El workflow \codepath{refined_basin} reintegra celdas previamente clasificadas como \codepath{unknown}. Compara cada trayectoria contra referencias target positiva y negativa mediante un puntaje de geometría:
\begin{equation}
S=\frac{w_c d_c+w_r d_r+w_f d_f+w_s d_s}{w_c+w_r+w_f+w_s},
\end{equation}
donde \(d_c\) es distancia de nube de cola, \(d_r\) distancia relativa de rangos, \(d_f\) diferencia relativa de FFT y \(d_s\) distancia de sección. Con los valores por defecto, se acepta target sólo si se satisfacen simultáneamente umbrales como:
\begin{equation}
S\leq 0.95,
\qquad d_c\leq0.85,
\qquad d_r\leq1.50,
\qquad d_f\leq1.00,
\qquad d_s\leq1.20.
\end{equation}
Estos umbrales son contratos operativos de clasificación, no teoremas.

\section{Análisis espectral, Lyapunov, bifurcaciones y Matignon}
\label{sec:analisis_adicionales_disponibles}

\subsection{FFT dominante y entropía espectral}

La API activa calcula FFT sobre la cola de la trayectoria. Para una serie \(x_n\), primero se remueve la media y se aplica ventana de Hann:
\begin{equation}
\tilde x_n=(x_n-\bar x)w_n.
\end{equation}
Luego
\begin{equation}
S_k=|\operatorname{rFFT}(\tilde x)_k|^2.
\end{equation}
La frecuencia dominante es
\begin{equation}
f_{\max}=f_{k^*},
\qquad k^*=\arg\max_{k\geq1}S_k.
\end{equation}
La entropía espectral normalizada es
\begin{equation}
H=-\frac{1}{\log N_f}\sum_{k\geq1}p_k\log(p_k+\varepsilon),
\qquad
p_k=\frac{S_k}{\sum_{j\geq1}S_j}.
\end{equation}
El campo correspondiente en los CSV suele aparecer como \codepath{fft_peak} o \codepath{fft_dominant_frequency}; la entropía aparece como \codepath{psd_entropy}, aunque en la API activa puede provenir de la FFT con ventana y no necesariamente de Welch.

\subsection{PSD de Welch}

La PSD de Welch está conservada como análisis opcional en el pipeline legado. Se activa con:
\begin{lstlisting}
hidden-attractors-unified-chua --run-mode balanced --spectral --psd
\end{lstlisting}
En Linux/macOS:
\begin{lstlisting}
python -m hidden_attractors.workflows.unified_chua --run-mode balanced --spectral --psd
\end{lstlisting}
La lectura correcta es: FFT y PSD caracterizan contenido frecuencial y ayudan a comparar persistencia espectral, pero no prueban caos ni ocultedad.

\subsection{Exponentes de Lyapunov}

El repositorio conserva un backend C tipo Benettin para estimación operacional de exponentes de Lyapunov. Se activa desde el pipeline legado con:
\begin{lstlisting}
hidden-attractors-unified-chua --run-mode balanced --lyapunov
\end{lstlisting}
Para hacer que un fallo detenga la corrida:
\begin{lstlisting}
hidden-attractors-unified-chua --run-mode balanced --lyapunov --lyapunov-strict
\end{lstlisting}
La interpretación mínima es:
\begin{equation}
\lambda_{\max}>0
\quad\Rightarrow\quad
\text{evidencia operacional de sensibilidad caótica bajo el contrato probado.}
\end{equation}
No sustituye la prueba de ocultedad, porque la ocultedad es una propiedad de la cuenca respecto de los equilibrios, no sólo de la dinámica sobre el atractor.

\subsection{Diagramas de bifurcación}

La API activa contiene post-procesamiento para diagramas de bifurcación. No hace continuación de ramas en sentido estricto; toma trayectorias ya integradas para distintos parámetros y extrae máximos, mínimos o muestras de una observable. Matemáticamente, para cada parámetro \(\rho_i\), se conserva una colección
\begin{equation}
\mathcal B_i=\{x(t_k;\rho_i):t_k\geq t_{\rm burn},\ x(t_k)\text{ máximo local}\}.
\end{equation}
La función pública es:
\begin{lstlisting}
from hidden_attractors.analysis import bifurcation_points_from_trajectories
from hidden_attractors.plotting import plot_bifurcation_diagram
\end{lstlisting}
Modos disponibles:
\begin{itemize}
\item \codepath{maxima}: máximos locales;
\item \codepath{minima}: mínimos locales;
\item \codepath{both}: máximos y mínimos;
\item \codepath{sample}: muestra directa de la cola.
\end{itemize}

\subsection{Equilibrios y gráfica Matignon}

El script de análisis de equilibrios calcula para cada equilibrio:
\begin{itemize}
\item región: central, saturación izquierda, saturación derecha o frontera de conmutación;
\item autovalores del Jacobiano local;
\item margen fraccionario
\begin{equation}
\mu_{\rm Matignon}=\min_i\left(|\arg(\lambda_i)|-\frac{q\pi}{2}\right);
\end{equation}
\item etiqueta \codepath{matignon_stable}.
\end{itemize}
El archivo producido es \code{equilibria\_summary.csv}.
Una gráfica defendible debe mostrar, para cada equilibrio, si
\begin{equation}
\mu_{\rm Matignon}>0
\end{equation}
o no. Un ejemplo mínimo para generar una gráfica a partir del CSV es:
\begin{lstlisting}
import csv
import matplotlib.pyplot as plt

rows = list(csv.DictReader(open('equilibria_summary.csv', newline='')))
labels = [r['eq_id'] for r in rows]
margins = [float(r['min_arg_margin']) for r in rows]

plt.figure(figsize=(6,4))
plt.axhline(0.0, color='k', lw=0.8)
plt.bar(labels, margins)
plt.ylabel(r'$\min_i(|\arg\lambda_i|-q\pi/2)$')
plt.xlabel('equilibrio')
plt.tight_layout()
plt.savefig('matignon_equilibria_margin.png', dpi=220)
\end{lstlisting}
Esta figura no dice si el atractor es oculto; sólo clasifica estabilidad local de los equilibrios.

\section{Inventario de funciones disponibles, activas y opcionales}
\label{sec:inventario_funciones_disponibles}

El repositorio contiene una API activa en \codepath{hidden_attractors/} y scripts legacy todavía útiles en \codepath{tools/legacy/}. La siguiente tabla resume qué existe y cómo debe leerse.

\begin{center}
\small
\begin{tabular}{@{}L{0.28\textwidth}L{0.24\textwidth}L{0.38\textwidth}@{}}
\toprule
Bloque & Estado & Función disponible\\
\midrule
\codepath{models.chua} & API activa & Parámetros, campo vectorial y equilibrios del Chua no suave.\\
\codepath{analysis.trajectory} & API activa & Rangos, varianza, FFT, entropía espectral, secciones, distancias de nubes y casos de robustez.\\
\codepath{analysis.bifurcation} & API activa & Extracción de puntos para diagramas de bifurcación desde trayectorias ya calculadas.\\
\codepath{plotting.dynamics} & API activa & Retratos de fase, proyecciones, series temporales y diagramas de bifurcación.\\
\codepath{basins.classification} & API activa & Etiquetas \codepath{equilibrium}, \codepath{target_positive}, \codepath{target_negative}, \codepath{infinity}, \codepath{unknown}, \codepath{numerical_failure}.\\
\codepath{workflows.robustness_overlay} & CLI mantenida & Superposición de trayectorias al variar \(h,L_m,T\).\\
\codepath{workflows.sphere_controls} & CLI mantenida & Pruebas desde esferas alrededor de todos los equilibrios y robustez del candidato.\\
\codepath{workflows.refined_basin} & CLI mantenida & Reclasificación de celdas \codepath{unknown} mediante geometría contra referencias target.\\
\codepath{seed_generation} & API activa ligera & DF centrada, semilla sesgada y semillas Machado sin variables de entorno.\\
\codepath{solvers.FractionalHistory} & API activa ligera & Ventana de memoria EFORK; la integracion pesada se ejecuta con backend C.\\
\codepath{workflows.unified_chua} & CLI/API mantenida & Entrada publica sin \codepath{\$env}; delega calculos pesados al backend C.\\
\codepath{integrations.external_tools} & API activa opcional & Adaptadores para \codepath{nolds}, \codepath{antropy}, PyDSTool y referencias externas.\\
\codepath{tools/legacy/run_extended_search.py} & Legacy útil & Búsqueda centrada, sesgada, Machado, \(\rho_H\), continuación multiparamétrica y nubes de semillas.\\
\codepath{tools/legacy/unified_nyquist_hidden_pipeline.py} & Legacy amplio & Pipeline completo con modos \codepath{balanced}, \codepath{fast}, \codepath{basin_only}, \codepath{q_sweep}, \codepath{df_compare}, \codepath{machado_sweep}, \codepath{machado_sweep_fast}.\\
\codepath{tools/legacy/danca2017_chua_abm_replication.py} & Legacy de referencia & Réplica ABM de Danca y comparación contra candidatos del proyecto.\\
\codepath{tools/legacy/positive_x_basin_sweep.py} & Pendiente de migración & Barrido de cuenca en región \(x>0\), especialmente alrededor de \(E_+\).\\
\codepath{tools/legacy/chua_frac_lyapunov_efork_benettin.c} & Opcional nativo & Estimación operacional de exponentes de Lyapunov.\\
\bottomrule
\end{tabular}
\end{center}

\subsection{Funciones opcionales desactivadas por defecto}

Varias capacidades existen pero no siempre se ejecutan en una corrida normal:
\begin{center}
\small
\begin{tabular}{@{}L{0.30\textwidth}L{0.36\textwidth}L{0.24\textwidth}@{}}
\toprule
Capacidad & Activación típica & Salida esperada\\
\midrule
PSD de Welch & \codepath{hidden-attractors-unified-chua}\newline\codepath{--psd} & \codepath{unified_psd_summary.json}, figuras PSD.\\
Exportación TISEAN & \codepath{hidden-attractors-unified-chua}\newline\codepath{--tisean} & Series exportadas para análisis externo.\\
Lyapunov operativo & \codepath{hidden-attractors-unified-chua}\newline\codepath{--lyapunov} & Resumen de exponentes y convergencia.\\
Lyapunov estricto & \codepath{hidden-attractors-unified-chua}\newline\codepath{--lyapunov-strict} & Falla la corrida si el cálculo falla.\\
Planos extra de cuenca & \codepath{hidden-attractors-unified-chua}\newline\codepath{--basin-planes} & Cortes \(xy,xz,yz\).\\
Cuenca 3D o vistas extra & \codepath{pendiente C} & Figuras adicionales.\\
Ilustración de ocultedad & \codepath{hidden-attractors-unified-chua}\newline\codepath{--hidden-illustration} & Figuras de vecindades y target/no-target.\\
Sweep de \(q\) & \codepath{hidden-attractors-unified-chua}\newline\codepath{--run-mode q_sweep} & \codepath{q_order_sweep/}.\\
Comparación DF/Machado & \codepath{hidden-attractors-unified-chua}\newline\codepath{--run-mode df_compare} & \codepath{df_seed_comparison/}.\\
Sweep Machado completo & \codepath{hidden-attractors-unified-chua}\newline\codepath{--run-mode machado_sweep} & \codepath{machado_sweep/}.\\
Sweep Machado rápido & \codepath{hidden-attractors-unified-chua}\newline\codepath{--run-mode machado_sweep_fast} & \codepath{machado_sweep_fast/}.\\
\bottomrule
\end{tabular}
\end{center}

\section{Ruta completa recomendada actualizada}
\label{sec:ruta_completa_actualizada}

Para el caso Chua fraccionario no suave se recomienda ejecutar y documentar en este orden:

\begin{enumerate}
\item Fijar sistema, parámetros, \(q,h,L_m,T,t_{\rm burn}\) y versión de código.
\item Calcular equilibrios, Jacobianos, autovalores y margen de Matignon.
\item Generar gráfica de márgenes de Matignon por equilibrio.
\item Ejecutar \codepath{q_sweep} para detectar órdenes \(q\) prometedores.
\item Ejecutar \codepath{df_compare} para comparar semilla clásica y Machado.
\item Ejecutar \codepath{machado_sweep_fast}; si hay candidatos, ejecutar \codepath{machado_sweep} completo.
\item Ejecutar \codepath{run_extended_search.py} para combinar centradas, sesgadas, Machado, \(\rho_H\), continuación y nubes de semillas.
\item Filtrar candidatos por \codepath{residual_abs}, \(\rho_H\), no divergencia, no convergencia a equilibrio y geometría acotada no trivial.
\item Reintegrar candidatos seleccionados con tiempo largo.
\item Calcular FFT, entropía espectral y, si se activa, PSD de Welch.
\item Calcular exponentes de Lyapunov si se requiere evidencia adicional de caos.
\item Construir diagramas de bifurcación para parámetros relevantes \(q,\alpha,\beta,\gamma,m_0,m_1\) si hay trayectorias por barrido.
\item Ejecutar \codepath{robustness_overlay} para \(h,L_m,T\).
\item Ejecutar \codepath{sphere_controls} para todos los equilibrios, radios y muestras.
\item Generar cuencas \(xy\), y si está habilitado, \(xz\), \(yz\) o cortes locales.
\item Ejecutar \codepath{refined_basin} si hay muchas celdas \codepath{unknown}.
\item Comparar todo contra la réplica ABM de Danca, separando diferencias de método: ABM con historia completa frente a EFORK con memoria truncada.
\item Redactar conclusión con nivel de evidencia, sin promover a \codepath{hidden_verified} salvo que no haya contacto target desde vecindades de todos los equilibrios bajo los contratos probados.
\end{enumerate}


\begin{thebibliography}{99}

\bibitem{Danca2017}
M. F. Danca.
\emph{Hidden chaotic attractors in fractional-order systems}.
Nonlinear Dynamics, 89(1), 2017, 577.
DOI relacionado: \texttt{10.1007/s11071-017-3472-7}.
Preprint: \texttt{arXiv:1804.10769}.

\bibitem{Caputo1967}
M. Caputo.
\emph{Linear Models of Dissipation whose Q is almost Frequency Independent-II}.
Geophysical Journal International, 1967.

\bibitem{Matignon1996}
D. Matignon.
\emph{Stability Results for Fractional Differential Equations with Applications to Control Processing}.
Computational Engineering in Systems Applications, 1996.

\bibitem{DiethelmFordFreed2002}
K. Diethelm, N. J. Ford, and A. D. Freed.
\emph{A Predictor-Corrector Approach for the Numerical Solution of Fractional Differential Equations}.
Nonlinear Dynamics, 2002.

\bibitem{LeonovKuznetsov2013}
G. A. Leonov and N. V. Kuznetsov.
\emph{Hidden Attractors in Dynamical Systems: From Hidden Oscillations in Hilbert-Kolmogorov, Aizerman, and Kalman Problems to Hidden Chaotic Attractor in Chua Circuits}.
International Journal of Bifurcation and Chaos, 2013.

\bibitem{TavazoeiHaeri2009}
M. S. Tavazoei and M. Haeri.
\emph{A Proof for Non Existence of Periodic Solutions in Time Invariant Fractional Order Systems}.
Automatica, 2009.

\bibitem{Machado2015}
J. Tenreiro Machado.
\emph{Fractional order describing functions}.
Signal Processing, 2015.

\end{thebibliography}

\end{document}


